%! Composition of Functions — Unit 6, Lesson 6.6 — Exit Ticket (Answer Key)
\documentclass[10pt]{article}
\usepackage{algebra2-article}
\usepackage{algebra2-key}   % pulls in -boxes; do NOT also load -boxes
\newcommand{\TallMath}[1]{$\displaystyle #1\rule[-1.4em]{0pt}{3.2em}$}
% In-formula answer slot (\ans is text-mode and must never appear inside $...$).
\newcommand{\mans}[1]{{\color{keyred}\mathbf{#1}}}

\begin{document}

\pageheader{Unit 6, Lesson 6.6}{Exit Ticket --- Answer Key}
\namedateperiod

\vspace{0.08in}

No notes. Work \textbf{inside out} --- and show the inside step, not just the final number.

\begin{enumerate}[leftmargin=1.8em, itemsep=10pt]
  \item \textbf{Both orders.} Let $f(x)=3x-2$ and $g(x)=x^{2}+1$.

        {\small \emph{Work:} (a) $g(2)=2^{2}+1=5$, then $f(5)=3(5)-2=13$. \\
        (b) $f(2)=3(2)-2=4$, then $g(4)=4^{2}+1=17$.}

        (a) $(f\circ g)(2)=\ans{$13$}$ \hspace{0.35in}
        (b) $(g\circ f)(2)=\ans{$17$}$

        \vspace{4pt}
        (c) Your two answers are different. Which function did you evaluate \emph{first} in part (a)?
        \ans{$g$ --- the inside function, the one written next to the $x$}

  \item \textbf{Algebra and domain.} Let $f(x)=\sqrt{x}$ and $g(x)=x+7$.

        {\small \emph{Work:} $(f\circ g)(x)=f(g(x))=f(x+7)=\sqrt{x+7}$. \; Stage 1 ($g$) accepts
        every real number; stage 2 ($f$) accepts only inputs $\ge 0$, so we need $x+7\ge0$.}

        (a) $(f\circ g)(x)=\ans{$\sqrt{x+7}$}$ \hspace{0.3in}
        (b) Domain: \ans{$x\ge-7$, i.e.\ $[-7,\infty)$}

        \vspace{5pt}
        (c) In one sentence, say why the domain is \emph{not} all real numbers --- your sentence
        should name which \emph{stage} of the composition does the refusing. \ansline{Stage 1 takes
        anything, but stage 2 is a square root, so the number handed to it ($x+7$) must be
        non-negative --- which forces $x\ge-7$.}

  \item \textbf{SOL-style multiple choice.} If $f(x)=x^{2}$ and $g(x)=x-5$, which expression is
        equal to $(f\circ g)(x)$?
        \begin{enumerate}[label=\Alph*., leftmargin=1.9em, itemsep=1pt]
          \item $x^{2}-5$
          \item $(x-5)^{2}$ \quad \textcolor{keyred}{\textbf{$\leftarrow$ correct}}
          \item $x^{2}-25$
          \item $x^{3}-5x^{2}$
        \end{enumerate}
        \vspace{2pt}
        \begin{itemize}[leftmargin=1.4em, nosep]
          \item \textbf{B} is correct: $g$ runs first, so $f$ receives $x-5$ and squares
                \emph{all} of it.
          \item \textbf{A} is $(g\circ f)(x)$ --- the \textbf{order error}, and the most common miss
                on this item. Square first, subtract second.
          \item \textbf{C} expands $(x-5)^{2}$ as $x^{2}-25$: the perfect-square trinomial error from
                Unit 4, still alive. Correct expansion is $x^{2}-10x+25$.
          \item \textbf{D} \emph{multiplied} the two functions instead of composing them. A student
                who picks this thinks $\circ$ is a product symbol.
        \end{itemize}
\end{enumerate}

\begin{teachernote}
\small Graded for completion. Sort into four piles; they decide how 6.7 opens.
\begin{itemize}[leftmargin=1.3em, nosep]
  \item \textbf{Got it} --- $13$ and $17$, $\sqrt{x+7}$ with $x\ge-7$, item 3 $=$ B.
  \item \textbf{Order reversed} --- item 1 answered $17$ and $13$, and/or item 3 $=$ A. This is the
        signature error of the lesson and it is pure notation: the function written next to the $x$
        goes first. Re-run Warm-Up 1(d) with them; they got it right before the symbol existed.
  \item \textbf{Composed correctly, domain missing} --- item 2(b) answered ``all reals.'' They are
        simplifying and then reading the domain off the answer. The two-stage gate needs re-teaching
        \emph{before} 6.7, because the entire restricted-domain argument there depends on it.
  \item \textbf{Composition $=$ multiplication} --- item 3 $=$ D. Rare but serious; sit with them and
        run one composition as a two-box machine diagram.
\end{itemize}
\textbf{Item 2(c) is the conceptual item.} A student who writes ``because of the square root'' has
half of it; the sentence you want names \emph{which stage} refuses and \emph{what number} it is
refusing --- not $x$, but $x+7$. That distinction is the whole content of the domain rule, and it is
what makes the $\left(\sqrt{x}\,\right)^{2}=x$ surprise land in Lesson 6.7.
\end{teachernote}

\end{document}
